\section{Experimental Setup}

We designed experiments to test our central claim --- that three-layer contract-based validation eliminates downstream failures from feasibility constraint violations --- through a hierarchical hypothesis decomposition. Our experimental approach validates three causal mechanisms before testing the end-to-end benefit, ensuring we understand not just whether the approach works but why it works.

\subsection{Hypotheses and Experimental Questions}

We decompose the main hypothesis into four testable sub-hypotheses:

\textbf{h-e1 (EXISTENCE):} Can three-layer contract-based validation be implemented with acceptable overhead? Tests whether the validation framework is implementable and practical. Success criterion: >50\% detection improvement over schema-only with <100ms execution overhead and <50 LOC implementation.

\textbf{h-m1 (MECHANISM):} Do contracts force explicit checking of constraints schema cannot express? Tests the expressiveness gap between schema-only and contract-based validation. Success criterion: Contracts check 100\% of constraints, schema checks <50\%, gap $\geq$75 percentage points.

\textbf{h-m2 (MECHANISM):} Does three-layer validation detect violations single-layer misses? Tests whether layers are complementary (cumulative detection) rather than redundant. Success criterion: Three-layer detection exceeds schema-only by $\geq$40 percentage points, with pattern layer achieving $\geq$80\% recall on semantic constraints.

\textbf{h-m3 (MECHANISM + PRIMARY PREDICTION):} Does early detection prevent downstream failures? Tests end-to-end benefit of boundary validation. Success criterion: Failure rate reduction $\geq$80\% compared to schema-only baseline (primary prediction P1).

This decomposition tests causal chain steps sequentially: h-e1 validates feasibility $\rightarrow$ h-m1 validates expressiveness $\rightarrow$ h-m2 validates multi-layer detection $\rightarrow$ h-m3 validates failure prevention.

\subsection{Datasets and Test Corpora}

We use two complementary test datasets:

\textbf{Adversarial Test Suite (h-m2):} 25 hand-crafted test cases with known ground-truth violations, designed to measure detection rates per layer. Distribution: 10 structural violations (schema baseline testing), plus semantic/compositional violations across constraint types (C1 synthetic: 5 cases, C2 human eval: 5 cases, C3 standard dataset: 2 cases, C4 new benchmark: 3 cases). Each test case labeled with expected detection layer (schema, pattern, or contract). Purpose: Measure per-layer detection rates with controlled violation types.

\textbf{Placeholder Hypothesis Corpus (h-m3):} 100 minimal hypotheses with typed interfaces satisfying schema contracts but minimal substantive content. 80 valid hypotheses (no violations), 20 with embedded constraint violations (5 per constraint type C1--C4). Purpose: Test end-to-end pipeline execution under schema-only versus contract-based validation conditions, measuring Phase 4/5 failure rates.

\textbf{Rationale for placeholder content:} Infrastructure testing requires minimal artifact fidelity, analogous to drone Software-in-the-Loop (SIL) testing where simulated sensor inputs test control logic without physical hardware. Placeholder hypotheses satisfy typed phase interfaces, which is sufficient to test constraint enforcement logic without requiring substantive research content. This approach enables rapid iteration on validation framework design.

\subsection{Baselines}

\textbf{Schema-only validation (primary baseline):} Pydantic BaseModel with typed fields, Field constraints (min\_length, Literal enums), but no field validators or model validators. This baseline isolates pure structural schema expressiveness --- what can be validated through type checking and structural constraints alone. Field validators are excluded by design choice to measure the expressiveness gap, not because they are unavailable in Pydantic. This creates a conservative baseline (including validators would increase baseline performance and reduce measured gaps).

\textbf{No validation (control):} Direct phase execution without boundary validation. Establishes upper bound on failure rates when constraints are not enforced.

\textbf{Rationale:} We compare against structural-schema-only rather than manual review because automated pipelines require automated validation. Manual review provides high accuracy but does not scale to multi-agent research automation. We exclude field validators from the baseline to isolate what pure schema validation (types, structure) can express versus what requires explicit semantic/compositional checking (the research question).

\subsection{Evaluation Metrics}

\textbf{Expressiveness Gap (h-m1):} Percentage point difference in constraint coverage between contracts and schema-only. Computed as: (contracts\_coverage - schema\_coverage) $\times$ 100. Measures whether contracts enable checking schema cannot perform.

\textbf{Detection Rate (h-m2):} Percentage of ground-truth violations detected at phase boundaries. Computed per layer (schema, pattern, contract) and cumulatively (all three layers). Measures multi-layer validation effectiveness.

\textbf{Failure Reduction (h-m3, P1):} Percentage reduction in Phase 4/5 downstream failures. Computed as: (baseline\_failures - contract\_failures) / baseline\_failures $\times$ 100. Primary prediction P1 requires $\geq$80\% reduction.

\textbf{Implementation Cost (h-e1):} Lines of code for contract layer specification, execution time overhead per validation call. Measures practical feasibility.

\subsection{Experimental Procedure}

\textbf{For h-e1 (Existence):} Implement three-layer validation for Phase 2A$\rightarrow$2B boundary. Run adversarial test suite, measure: (1) detection rate improvement vs schema-only, (2) execution overhead, (3) contract layer LOC, (4) false positive rate on valid inputs.

\textbf{For h-m1 (Expressiveness):} Analyze constraint definitions (C1--C4) to determine which are expressible via schema-only versus contracts. For each constraint, attempt to encode in schema alone (Pydantic Field, Literal, validators disabled). Record coverage: constraints\_expressible / total\_constraints. Repeat for contract-based (all layers enabled).

\textbf{For h-m2 (Multi-layer detection):} Run adversarial test suite through three conditions: (1) schema-only, (2) schema + pattern, (3) schema + pattern + contract. For each violation, record which layer detects it. Compute cumulative detection rates and per-layer contributions.

\textbf{For h-m3 (Failure prevention):} Run 100-case placeholder corpus under two conditions: (A) schema-only validation at boundaries, (B) contract-based validation at boundaries. For each hypothesis, track: (1) violations detected at boundaries (Phase 2$\rightarrow$3$\rightarrow$4), (2) failures at Phase 4/5 implementation. Compute failure rate per condition and reduction.

\subsection{Threats to Validity}

\textbf{Internal validity:} Schema-only baseline excludes field validators (pattern layer) to isolate schema expressiveness. This design choice is conservative --- including validators in the baseline would increase baseline performance, reducing measured gaps. However, our research question is ``what can schema express'' not ``what can Pydantic express,'' making validator exclusion appropriate.

\textbf{External validity:} Placeholder content has low fidelity relative to substantive research. Results demonstrate that constraint enforcement logic functions correctly on typed interfaces but do not prove the approach generalizes to real research content with complex semantic dependencies. This limitation is acknowledged and deferred to future work (Phase 5 baseline comparison on real pipeline executions).

\textbf{Construct validity:} Detection rate measures violations caught, but not false positives (valid inputs rejected). We measure false positive rate on valid test cases to ensure contracts do not over-restrict. Zero false positives observed in validation.

\textbf{Statistical power:} Deterministic test suite (known ground truth violations) provides perfect reproducibility but no statistical inference. Failure rates are measured counts, not probabilistic samples, so significance testing is not applicable. Effect sizes (100pp expressiveness gap, 48pp detection gap, 100\% failure reduction) are large enough that statistical uncertainty is negligible.
